\chapter{Calendar spreads et trades de courbe}\label{ch:s1:calendar-spreads-and-curve-trades}

Le 20 avril 2020, le contrat de mai sur le pétrole brut WTI a ouvert à 17.73~\$ le baril et a clôturé à moins 37.63~\$, le premier prix négatif en 37 ans d'existence~;
les services de la CFTC ont constaté un marché en surproduction, un effondrement de la demande et un point de livraison, Cushing dans l'Oklahoma, dont le stockage utile était proche de sa capacité
depuis mars. Un calendar spread, achat d'une échéance et vente de la suivante, paraît le trade le plus sûr parmi les contrats à terme~: les deux jambes partagent presque tout leur risque
et le spread évolue lentement. Sur les trente-neuf ans de cours de compensation du WTI publiés par l'EIA, une règle qui mise sur le \omterm{def:s1:calendar-spreads-and-curve-trades:mr}{retour à la moyenne du spread} a obtenu un ratio de Sharpe de
0.30 avant coûts et 0.05 après, de 1986 à 2019~; en 2020 elle a perdu 43~\% de son capital pour une cible de volatilité de 10~\%, surtout en mars, alors que le stockage
se remplissait. La construction porte le nom \texttt{firm.curvestrat}.

\section{Trader la forme d'une courbe}

Une courbe de contrats à terme a un niveau, une pente et, plus loin, une courbure. Une position outright porte sur le niveau~; un calendar spread (livre~1, chapitre~19)
porte sur la pente entre deux échéances et est presque insensible au niveau.

\begin{definition}[Trade de courbe]\label{def:s1:calendar-spreads-and-curve-trades:curve}
Un \emph{trade de courbe}\index{trade de courbe} (curve trade) est une position sur deux échéances ou plus d'un même contrat à terme (ou sur des points d'une courbe de rendement ou de volatilité)
dont les pondérations annulent l'exposition au niveau de la courbe, de sorte que son P\&L résulte des variations de la forme de la courbe~: un calendar spread pour la pente entre
deux échéances, un butterfly pour la courbure entre trois.
\end{definition}

\begin{proposition}[Un calendar spread porte sur la variation du carry]\label{prop:s1:calendar-spreads-and-curve-trades:carry}
Si le logarithme du prix à terme pour l'échéance $T$ est $\ln F_t(T) = x_t - c_t (T - t) + g(T)$, avec $x_t$ un niveau commun, $c_t$ le carry et $g$ une
composante saisonnière fixée par la date de livraison, alors une position acheteuse sur l'échéance $T_a$ et vendeuse sur l'échéance $T_b > T_a$ a un P\&L logarithmique de
\[
d\big(\ln F_t(T_a) - \ln F_t(T_b)\big) = (T_b - T_a)\, dc_t .
\]
\end{proposition}
\begin{proof}
Le niveau $x_t$ et les termes saisonniers s'annulent ou sont constants~; $d[-c_t(T_a - t) + c_t(T_b - t)] = (T_b - T_a)\,dc_t$ puisque les termes en $c_t\,dt$
se compensent.
\end{proof}

Une position acheteuse sur le spread (proche moins lointain) gagne quand la courbe évolue vers la backwardation et perd quand le contango s'accentue~; sa taille croît avec la distance
entre les échéances. Ce qui fait bouger la pente des matières premières, c'est l'inventaire~: Gorton, Hayashi et Rouwenhorst ont montré que le rendement d'opportunité (convenience yield) diminue, de façon non linéaire,
à mesure que les inventaires augmentent, et que la base et les rendements passés reflètent l'état des inventaires et anticipent les primes de risque. Des inventaires faibles rendent la
courbe en backwardation et pentue~; des inventaires élevés la poussent en contango jusqu'au coût de stockage, le portage complet (full carry) du livre~3, chapitre~10.

\section{Retour à la moyenne des spreads}

\begin{definition}[Retour à la moyenne du spread]\label{def:s1:calendar-spreads-and-curve-trades:mr}
Le \emph{retour à la moyenne du spread}\index{retour à la moyenne du spread} (spread mean reversion) est la tendance d'un spread entre prix liés à revenir vers un niveau normal après
un écart~; une stratégie de retour à la moyenne du spread vend le spread lorsqu'il est haut par rapport à sa moyenne récente et l'achète lorsqu'il est bas.
\end{definition}

Trader un spread WTI sur plusieurs décennies exige de suivre l'identité des contrats. L'EIA publie des séries génériques~: le contrat~1 est celui qui arrive à échéance en premier, et le lendemain
d'une échéance, il devient l'ancien contrat~2. Le \Cref{lst:s1:calendar-spreads-and-curve-trades:pair} numérote les contrats par numéro de série et garde la paire la plus proche
à l'écart de toute échéance (1--2, passant à 2--3 cinq jours de bourse avant l'échéance du contrat proche), et suit les mêmes deux contrats d'un jour à l'autre, de sorte qu'un roll
de la série générique ne soit pas confondu avec une variation de prix. Le spread logarithmique $\ln(F_{\text{near}}/F_{\text{far}})$ de la paire détenue est persistant mais
qui revient à sa moyenne~: son autocorrélation journalière est de 0.949, soit une demi-vie de 13.3 jours de bourse.

La règle (\cref{lst:s1:calendar-spreads-and-curve-trades:band}) calcule le z-score du spread logarithmique par rapport à ses 20 derniers jours, prend position à contre-sens au-delà de 1.5
(achat du spread quand il est anormalement bas, vente quand il est haut), et sort quand le z-score revient à moins de 0.5 en valeur absolue ou change de signe. La position est
dimensionnée à l'entrée pour une volatilité annuelle de 10~\%, à partir d'une prévision EWMA du P\&L en dollars du spread, et chaque unité échangée coûte un cent par baril (hypothèse).

\begin{center}
\small
\begin{tabular}{@{}p{6.4cm}cccc@{}}
\toprule
filtre de stockage~: pas de spread acheteur quand le carry est inférieur à & aucun & $-80~\%$ & $-40~\%$ & $-20~\%$\\
\midrule
ratio de Sharpe avant coûts, 1986--2019 & 0.30 & 0.33 & 0.31 & 0.39\\
ratio de Sharpe après coûts, 1986--2019 & 0.05 & 0.08 & 0.05 & 0.14\\
part des jours avec une position & 47~\% & 47~\% & 45~\% & 41~\%\\
2020~: mars & $-58.3~\%$ & $-20.8~\%$ & $-10.5~\%$ & $-1.5~\%$\\
2020~: avril & $+22.9~\%$ & 0.0~\% & 0.0~\% & 0.0~\%\\
2020~: l'année & $-42.7~\%$ & $-28.1~\%$ & $-17.8~\%$ & $-8.8~\%$\\
\bottomrule
\end{tabular}
\end{center}

Le retour à la moyenne est réel mais modeste~: un ratio de Sharpe de 0.30 avant coûts sur 34 ans, et un cent par baril et par unité échangée en absorbe l'essentiel
(\cref{fig:s1:calendar-spreads-and-curve-trades:cumulative}). Par ailleurs, le bilan n'est pas stable~: après coûts, la règle affichait un ratio de Sharpe de 0.47 entre
1986 et 2004 et de $-0.32$ entre 2005 et 2019, si bien qu'elle perdait de l'argent depuis quinze ans avant que le stockage ne soit saturé. Le spread est un trade lent, d'apparence peu coûteux
dont l'edge tient à l'ampleur de ses coûts.

\begin{omfigure}
\begin{tikzpicture}
\begin{axis}[width=11cm,height=5.2cm,xlabel={\small année},ylabel={\small rendement cumulé (\%)},tick label style={font=\footnotesize},
  xmin=1986,xmax=2021,ymin=-40,ymax=90,grid=major,grid style={gray!25},legend columns=2,legend style={font=\scriptsize,
  at={(0.5,-0.3)},anchor=north,draw=none,fill=none,/tikz/every even column/.append style={column sep=8pt}},
  table/col sep=comma,/pgf/number format/1000 sep={}]
\addplot[omDef,thick] table[x=year,y=none]{figdata/strategies-1/21-calendar-spreads-and-curve-trades/cumulative.csv};
\addplot[omThm,thick] table[x=year,y=filter]{figdata/strategies-1/21-calendar-spreads-and-curve-trades/cumulative.csv};
\legend{sans filtre,pas de spread acheteur quand le carry est inférieur à $-20~\%$}
\end{axis}
\end{tikzpicture}

\omcaption{La règle de retour à la moyenne sur calendar spread WTI après coûts, 1986--2020, pour une cible de volatilité de 10~\%~: somme cumulée des rendements quotidiens, sans
filtre de stockage et avec le filtre qui s'avère le meilleur a posteriori. Données~: \texttt{s1\_curve.book}.}\label{fig:s1:calendar-spreads-and-curve-trades:cumulative}
\end{omfigure}

\section{Limites de stockage et prix négatif de 2020}

\begin{definition}[Risque de saturation des stocks]\label{def:s1:calendar-spreads-and-curve-trades:storage}
Le \emph{risque de saturation des stocks}\index{risque de saturation des stocks} (storage-limit risk) est le risque que le contango d'une matière première s'élargisse au-delà du coût de stockage lorsque les capacités de stockage sont épuisées, de sorte
qu'une position pariant sur le resserrement du spread (achat de l'échéance proche, vente de l'échéance lointaine) perde sans la borne que le portage complet
assure normalement.
\end{definition}

Le contango est normalement borné~: lorsque le contrat lointain coûte plus cher que le proche augmenté du coût de stockage du pétrole jusqu'à cette échéance, un trader qui dispose de capacités de stockage
achète le proche, le stocke et vend le lointain. Quand le stockage est plein, personne ne peut faire cette opération, et rien ne freine le spread. C'est là ce à quoi la règle s'est heurtée.
Le carry de la paire était de $+4.5~\%$ par an au début de janvier 2020~; la règle a pris une position acheteuse sur le spread le 10 janvier, dès qu'il est passé sous sa moyenne récente, et est restée
acheteuse, par intermittence, jusqu'au 23 avril (\cref{fig:s1:calendar-spreads-and-curve-trades:carry2020}). Au 31 mars, le carry était de $-216\%$ par an~; le 20 avril,
$-302~\%$~; le 21 avril, lendemain du jour où le contrat de mai a clôturé à $-\$37.63$, le carry de la paire juin--juillet était de $-575~\%$. La règle a perdu 58.3~\% du capital en
mars et en a regagné 22.9~\% en avril, car elle a repris des positions de taille réduite et le spread s'est partiellement redressé.

\begin{omfigure}
\begin{tikzpicture}
\begin{axis}[width=11cm,height=5.2cm,xlabel={\small jour de négociation de 2020},ylabel={\small carry de la paire détenue (\% par an)},
  tick label style={font=\footnotesize},xmin=1,xmax=124,ymin=-600,ymax=50,grid=major,grid style={gray!25},
  xtick={1,41,63,84,104},xticklabels={2 janv.,2 mars,1 avr.,1 mai,1 juin},legend style={font=\scriptsize,at={(0.5,-0.3)},anchor=north},
  unbounded coords=jump,table/col sep=comma]
\addplot[omDef,thick] table[x=n,y=carry_pct]{figdata/strategies-1/21-calendar-spreads-and-curve-trades/carry2020.csv};
\addplot[only marks,mark=*,mark size=1.1pt,omThm] table[x=n,y=long]{figdata/strategies-1/21-calendar-spreads-and-curve-trades/carry2020.csv};
\addplot[black!50,dashed] coordinates {(1,-20) (124,-20)};
\legend{carry de la paire détenue,jours en position acheteuse sur le spread,filtre à $-20~\%$}
\end{axis}
\end{tikzpicture}

\omcaption{Le carry annualisé, $12\ln(F_{\text{near}}/F_{\text{far}})$, de la paire WTI détenue par la règle au premier semestre 2020, issu des cours de règlement de l'EIA,
et les jours où la règle non filtrée était en position acheteuse sur le spread. Données~: \texttt{s1\_curve.wti}.}\label{fig:s1:calendar-spreads-and-curve-trades:carry2020}
\end{omfigure}

Un filtre qui interdit et clôture les positions acheteuses sur spread en contango très prononcé réduit la perte~: à 28.1~\% avec un seuil de $-80~\%$ par an, 17.8~\% avec $-40~\%$,
8.8~\% avec $-20~\%$, qui présente aussi le meilleur ratio de Sharpe avant 2020 (0.14 après coûts). Il serait facile de présenter le filtre à $-20~\%$ comme la
stratégie. Il a été retenu après examen de 2020, parmi quatre candidats, sur un échantillon ne comportant qu'une seule crise de stockage. Avant 2020, le carry était inférieur à $-20~\%$ sur
13.1~\% des jours et inférieur à $-40~\%$ dans 3.5~\% d'entre eux. La version défendable du filtre découle du mécanisme et non du backtest~: mesurer à quel point
le stockage approche de la saturation (inventaires, coût de stockage) et cesser de parier contre le contango lorsque c'est le cas.

\section{Trades de courbe dans l'univers synthétique, taux et volatilité}

Sur les dix matières premières de \texttt{firm.synthfut}, la proposition est vérifiée exactement~: le spread logarithmique entre le contrat le plus proche et le contrat à douze mois se déplace
un pour un avec le carry injecté (une corrélation de 1.00), parce que les courbes du modèle n'ont pas de bruit propre. Le carry revient à la moyenne propre à chaque marché
avec une demi-vie d'un an, si bien qu'un portefeuille de spreads indexé sur le z-score du spread à 252 jours affiche un ratio de Sharpe de 1.00 avant coûts lorsqu'il est négocié chaque jour~; après 4
points de base par jambe et par unité échangée, il perd ($-0.29$). Rééquilibré chaque mois, il conserve 0.76 de ses 1.05 avant coûts. Un signal lent négocié chaque jour paie le
prix du bruit dans sa propre position.

La même structure apparaît dans d'autres courbes. Sur les taux d'intérêt, les \omterm{def:s1:calendar-spreads-and-curve-trades:curve}{trades de courbe} sont des steepeners et des flatteners entre deux échéances, ou des butterflies,
pondérés pour être neutres en duration (livre~2)~; leur P\&L est la variation de la pente de la courbe, comme ici. En volatilité, les contrats à terme sur le VIX sont généralement en contango,
car elles incorporent une prime sur le VIX spot attendu. Simon et Campasano ont montré que, de 2006 à 2011, la base des contrats à terme sur le VIX ne permettait pas de prévoir
les variations du VIX spot, mais permettait de prévoir les variations des cours des contrats à terme sur le VIX, et que vendre à découvert des contrats à terme sur le VIX en contango et les acheter en backwardation, avec couverture
par des contrats à terme sur mini-S\&P 500, a été très rentable et robuste aux coûts. Le risque de saturation de la courbe de volatilité, c'est le krach~: le contango qui rémunère la position vendeuse
est précisément ce qui s'inverse quand la volatilité bondit.

\section{Fiches de stratégie}

\begin{strategyfile}{Retour à la moyenne du spread front--second}\label{strat:s1:calendar-spreads-and-curve-trades:front}
\sfield{Qui paie, et pourquoi} Les hedgers et les roll d'indices qui éloignent le spread proche de son niveau normal pendant quelques semaines.
\sfield{Instruments et plateformes} Des calendar spreads sur les contrats à terme liquides de matières premières (WTI, Brent, produits raffinés, céréales), négociés comme spreads cotés.
\sfield{Signal} Le z-score du spread logarithmique par rapport à sa moyenne récente (20 jours).
\sfield{Dimensionnement et exécution} Entrer au-delà de 1.5, sortir en deçà de 0.5~; position dimensionnée à l'entrée selon une cible de volatilité~; paire éloignée de son échéance.
\sfield{Coûts} La fourchette cotée, d'un tick ou davantage~; elle est de la taille de l'edge.
\sfield{Comment elle s'éteint} \omterm{def:s1:calendar-spreads-and-curve-trades:storage}{Risque de saturation des stocks}, lorsque le contango s'élargit sans limite~; coûts.
\sfield{Horizon, capacité, infrastructure} Quelques semaines~; la capacité sur les spreads proches est importante~; données d'identité des contrats.
\sfield{Backtest honnête} Identité des contrats à travers les échéances~; cotations de spreads et coûts~; les crises de stockage présentes dans l'échantillon.
\sfield{Sources} CFTC interim staff report (2020)~; ce chapitre~: 0.30 avant coûts et 0.05 après, 1986--2019~; $-42.7~\%$ en 2020.
\end{strategyfile}

\begin{strategyfile}{Steepener de courbe sur contrats à terme sur matières premières}\label{strat:s1:calendar-spreads-and-curve-trades:steepener}
\sfield{Qui paie, et pourquoi} Le cycle des inventaires~: la pente revient à sa moyenne à mesure que les stocks sont puisés ou reconstitués.
\sfield{Instruments et plateformes} Contrats à terme sur matières premières, à échéance proche et à douze mois.
\sfield{Signal} La pente par rapport à sa moyenne sur un an~; données d'inventaires.
\sfield{Dimensionnement et exécution} Se placer acheteur du spread lorsque le contango est anormalement profond et que les inventaires baissent, vendeur lorsque la backwardation est anormalement forte~;
rééquilibré chaque mois.
\sfield{Coûts} Deux jambes, deux roll.
\sfield{Comment elle s'éteint} Inventaires qui continuent de gonfler~; \omterm{def:s1:calendar-spreads-and-curve-trades:storage}{risques de saturation des stocks}.
\sfield{Horizon, capacité, infrastructure} Mois.
\sfield{Backtest honnête} Données d'inventaires telles que publiées, avec leurs révisions~; la saisonnalité de la courbe séparée de sa pente.
\sfield{Sources} Gorton, Hayashi et Rouwenhorst (2013)~; ce chapitre~: 0.76 sur l'univers synthétique, mensuel.
\end{strategyfile}

\begin{strategyfile}{Calendar spread sur la période de roll}\label{strat:s1:calendar-spreads-and-curve-trades:roll}
\sfield{Qui paie, et pourquoi} Les \omterm{def:s1:index-rebalancing:fund}{fonds indiciels} et les ETF qui effectuent un roll de grosses positions selon des calendriers publiés, en pesant sur le spread concerné.
\sfield{Instruments et plateformes} Les spreads proches des matières premières indicielles pendant les fenêtres de roll.
\sfield{Signal} Le calendrier des rolls d'indices (les flux du chapitre~9 appliqués aux contrats à terme).
\sfield{Dimensionnement et exécution} Prendre la contrepartie du roll avant qu'il ait lieu, effectuer le débouclage pendant celui-ci.
\sfield{Coûts} Faibles~; concurrence de la part d'autres front-runners.
\sfield{Comment elle s'éteint} Les \omterm{def:s1:index-rebalancing:fund}{fonds indiciels} étalant leurs roll sur davantage de jours ou de dates.
\sfield{Horizon, capacité, infrastructure} Jours.
\sfield{Backtest honnête} Les règles des indices et les dates de roll en vigueur à chaque date.
\sfield{Sources} Aucun chiffre de performance vérifié.
\end{strategyfile}

\begin{strategyfile}{Trade de courbe sur contrats à terme de volatilité}\label{strat:s1:calendar-spreads-and-curve-trades:vix}
\sfield{Qui paie, et pourquoi} Les acheteurs de protection contre la volatilité, qui paient une prime au-dessus du VIX attendu.
\sfield{Instruments et plateformes} Contrats à terme sur le VIX, couverts par des contrats à terme sur indices actions.
\sfield{Signal} La base~: contrats à terme au-dessus du VIX spot (contango) ou en dessous (backwardation).
\sfield{Dimensionnement et exécution} Vendeur en contango, acheteur en backwardation, exposition actions couverte~; position réduite face au risque de krach.
\sfield{Coûts} Modérés~; roll.
\sfield{Comment elle s'éteint} Pics de volatilité qui inversent le contango en une journée.
\sfield{Horizon, capacité, infrastructure} Semaines.
\sfield{Backtest honnête} Cours de règlement des contrats à terme, et non de l'indice spot~; ratios de couverture estimés hors échantillon~; pics présents dans l'échantillon.
\sfield{Sources} Simon et Campasano (2014).
\end{strategyfile}

\section{Tutoriel~: moins trente-sept}\label{tut:s1:calendar-spreads-and-curve-trades:tut}

\begin{tutorial}
\textbf{Objectif.} Construire le calendar spread WTI par identité des contrats à partir des séries génériques de l'EIA, et négocier son retour à la moyenne avec et sans
filtre de stockage, suivre son évolution en 2020 et répéter l'opération sur les courbes synthétiques. \textbf{État final~:} le tableau et les deux figures.
\begin{tutsteps}
\item \textbf{Identité des contrats}~: numéros de série, la paire détenue et son P\&L à travers les échéances.
\omcode{code/firm/curvestrat/firm_curvestrat.py}{26}{59}{P\&L du calendar spread par identité des contrats.}\label{lst:s1:calendar-spreads-and-curve-trades:pair}
\item \textbf{La règle}~: entrer au-delà d'un z-score, sortir en deçà d'un z-score plus petit, avec un filtre sur les spreads acheteurs.
\omcode{code/firm/curvestrat/firm_curvestrat.py}{84}{97}{La règle de bande avec filtre de stockage.}\label{lst:s1:calendar-spreads-and-curve-trades:band}
\item \textbf{Lancer} \texttt{persistence()}, \texttt{book(threshold)} pour les quatre filtres, \texttt{synthetic(252, monthly)} dans les deux sens, et
\texttt{fig\_curve.py}.
\end{tutsteps}
\textbf{Pour aller plus loin.} Remplacer le filtre de carry par les données d'inventaires de l'EIA à Cushing~; ajouter un butterfly sur les contrats 1, 2 et 3~; facturer la
fourchette affichée du spread coté plutôt qu'un centime fixe.
\end{tutorial}

\section{Construction~: calendar spreads}\label{bld:s1:calendar-spreads-and-curve-trades:curvestrat}

\begin{build}
\bfield{Rôle} Calendar spreads par identité des contrats, leur P\&L à travers les échéances, signaux de retour à la moyenne et filtre de régime.
\bfield{Interface} \texttt{serials(expiry)}, \texttt{pair\_pnl(C, m, near, days\_before, expiry)}, \texttt{zscore(x, window)},
\texttt{positions(z, vol, target, cap, block)}, \texttt{band(z, vol, enter, leave, target, block)}, \texttt{spread\_carry(log\_spread, days)}.
\bfield{Règles} Un contrat conserve son numéro de série à travers les échéances~; la paire sort avant les derniers jours du contrat proche~; les signaux utilisent les données jusqu'à $t$.
\bfield{Tests d'acceptation} \texttt{code/firm/curvestrat/tests/}~: une courbe construite à la main traversant une échéance, avec le P\&L de la paire calculé à la main~; le z-score, la
bande à l'entrée et à la sortie, le filtre~; le carry à partir d'un spread logarithmique.
\bfield{Pour aller plus loin} Butterflies~; filtres d'inventaires~; cotations des spreads cotés.
\end{build}

\begin{omsources}
\item CFTC Division of Market Oversight and Office of the Chief Economist, \emph{Interim Staff Report on Trading in NYMEX WTI Crude Oil Futures Contract
Leading up to, on, and around April 20, 2020}, novembre 2020.
\item G.~B.~Gorton, F.~Hayashi et K.~G.~Rouwenhorst, «~The fundamentals of commodity futures returns~», \emph{Review of Finance} 17(1), 2013.
\item D.~P.~Simon et J.~Campasano, «~The VIX futures basis: evidence and trading strategies~», \emph{Journal of Derivatives} 21(3), 2014.
\item US Energy Information Administration, NYMEX WTI crude oil futures daily settlements, contrats 1 à 4.
\end{omsources}

\section{Exercices}

\begin{exercise}[$\star$]\label{exo:s1:calendar-spreads-and-curve-trades:1}
Un spread logarithmique a une autocorrélation quotidienne de 0.949. Quelle est sa demi-vie en jours de négociation~?
\end{exercise}

\begin{exercise}[$\star$]\label{exo:s1:calendar-spreads-and-curve-trades:2}
Deux contrats distants de 21 jours de négociation ont un spread logarithmique de $-0.01$. Quel est le carry annualisé~? Le 21 avril 2020, le carry de la paire détenue était de $-575.5~\%$ par
an~: quel était le rapport entre le prix du contrat proche et celui du contrat éloigné~?
\end{exercise}

\begin{exercise}[$\star$]\label{exo:s1:calendar-spreads-and-curve-trades:3}
Pourquoi la série générique «~contrat 1 moins contrat 2~» fausse-t-elle un backtest de spread le lendemain d'une échéance~?
\end{exercise}

\begin{exercise}[$\star\star$]\label{exo:s1:calendar-spreads-and-curve-trades:4}
La règle non filtrée affiche des ratios de Sharpe de 0.30 avant coûts et de 0.05 après, pour une volatilité d'environ 10~\%. Quel est, en ordre de grandeur, le montant que les coûts prélèvent chaque année~?
\end{exercise}

\begin{exercise}[$\star\star$]\label{exo:s1:calendar-spreads-and-curve-trades:5}
Pourquoi la règle a-t-elle perdu en mars 2020 et gagné en avril~?
\end{exercise}

\begin{exercise}[$\star\star$]\label{exo:s1:calendar-spreads-and-curve-trades:6}
Pourquoi le résultat du filtre de $-20~\%$ n'est-il pas une preuve de son efficacité~?
\end{exercise}

\begin{exercise}[$\star\star\star$]\label{exo:s1:calendar-spreads-and-curve-trades:7}
\emph{Programmation.} Exécuter \texttt{synthetic} à fréquences quotidienne et mensuelle. Pourquoi le même signal perd-il après coûts en fréquence quotidienne et affiche-t-il un ratio de Sharpe de 0.76 en fréquence mensuelle~?
\end{exercise}

\begin{exercise}[$\star\star\star$]\label{exo:s1:calendar-spreads-and-curve-trades:8}
\emph{Trouver la faille.} «~Le spread de front est revenu à sa moyenne chaque mois depuis vingt ans~; on peut le faire tourner avec un effet de levier dix fois supérieur, car les jambes se
couvrent mutuellement.~»
\end{exercise}

\section{Problème~: Moins trente-sept}

\begin{problem}[{Problème du week-end --- un spread rencontre une limite de stockage}]\label{pb:s1:calendar-spreads-and-curve-trades:1}
Les règlements-livraisons du WTI publiés par l'EIA, les courbes synthétiques et le dossier public.

\medskip\noindent\textbf{Partie I --- Courbes.}
\begin{enumerate}
\item Définir un \omterm{def:s1:calendar-spreads-and-curve-trades:curve}{trade de courbe} et le \omterm{def:s1:calendar-spreads-and-curve-trades:mr}{retour à la moyenne du spread}.
\item Énoncer et démontrer la proposition sur les calendar spreads et le carry.
\item Que Gorton, Hayashi et Rouwenhorst ont-ils établi au sujet des inventaires~?
\item Qu'est-ce qui borne le contango, et quand cette borne échoue-t-elle~?
\end{enumerate}

\medskip\noindent\textbf{Partie II --- Le spread du WTI.}
\begin{enumerate}[resume]
\item Comment suit-on le spread selon l'identité des contrats~?
\item Quelle est la persistance du spread logarithmique~?
\item Décrire la règle et ses coûts.
\item Donner ses ratios de Sharpe avant et après coûts, de 1986 à 2019.
\end{enumerate}

\medskip\noindent\textbf{Partie III --- 2020.}
\begin{enumerate}[resume]
\item Que constate le rapport des services de la CFTC au sujet du 20 avril 2020~?
\item Définir le \omterm{def:s1:calendar-spreads-and-curve-trades:storage}{risque de saturation des stocks}.
\item Décrire le parcours de la règle en mars et avril 2020.
\item Que fait chaque filtre, et pourquoi le meilleur est-il suspect~?
\end{enumerate}

\medskip\noindent\textbf{Partie IV --- Le verdict.}
\begin{enumerate}[resume]
\item Énoncer le \emph{résultat nommé}~: le ratio de Sharpe de la stratégie de spread avant 2020 et sa perte d'avril 2020.
\item Que devrait utiliser un filtre défendable~?
\item Que montrent les courbes synthétiques quant à la vitesse de négociation~?
\item Que constatent Simon et Campasano au sujet de la base du contrat à terme sur le VIX~?
\item Quel est, pour la courbe de volatilité, l'équivalent d'une limite de stockage~?
\item Quelle fiche de stratégie dépend des règles des \omterm{def:s1:index-rebalancing:fund}{fonds indiciels}~?
\item Comment le carry de ce chapitre se rapporte-t-il à celui du chapitre~20~?
\item En une phrase~: que couvre un calendar spread, et que ne couvre-t-il pas~?
\end{enumerate}
\end{problem}

\section{Questions d'entretien}

\begin{interviewq}[$\star$ \iqroles{chercheur}]\label{iq:s1:calendar-spreads-and-curve-trades:1}
Qu'est-ce qui détermine la pente de la courbe des contrats à terme sur matières premières~?
\end{interviewq}

\begin{interviewq}[$\star\star$ \iqroles{développeur}]\label{iq:s1:calendar-spreads-and-curve-trades:2}
Comment stockeriez-vous les données de contrats à terme pour que les backtests de spreads suivent correctement les contrats à travers les échéances~?
\end{interviewq}

\begin{interviewq}[$\star\star$ \iqroles{trader}]\label{iq:s1:calendar-spreads-and-curve-trades:3}
Vous êtes long du spread de front du WTI et le contango s'élargit de 50~\% par an en deux jours. Que vérifiez-vous, et que faites-vous~?
\end{interviewq}

\begin{interviewq}[$\star\star$ \iqroles{risque}]\label{iq:s1:calendar-spreads-and-curve-trades:4}
Comment un book de calendar spreads doit-il être soumis aux appels de marge et aux tests de résistance~?
\end{interviewq}

\begin{interviewq}[$\star\star$ \iqroles{chercheur}]\label{iq:s1:calendar-spreads-and-curve-trades:5}
Pourquoi des positions vendeuses sur contrats à terme sur le VIX en contango pourraient-elles être rentables en moyenne, et quel est le risque~?
\end{interviewq}

\begin{interviewq}[$\star\star\star$ \iqroles{chercheur}]\label{iq:s1:calendar-spreads-and-curve-trades:6}
Un spread logarithmique suit un AR(1) de coefficient $\phi$ et d'écart-type d'innovation $\sigma$. Un trader détient $-\,s_t$ unités à la clôture de $t$.
Calculez l'espérance du P\&L quotidien et son ratio de Sharpe, et montrez comment un coût $k$ par unité échangée le modifie.
\end{interviewq}
